\chapter{规划策略层：基于任务分离的课程学习协同规划}

\section{规划层问题建模与框架}
上层制导系统负责多机协同与路径规划，输入全局观测与局部状态，输出期望的速度矢量或航点，决策频率为 10Hz。

\section{分阶段训练与课程学习}
构建从易到难的任务序列：
\begin{itemize}
    \item \textbf{Stage 1 (Single Agent)}: 单机悬停与导航，习得基础飞行动力学。
    \item \textbf{Stage 2 (Static Group)}: 多机编队与静态避障。
    \item \textbf{Stage 3 (Dynamic Pursuit)}: 动态目标围捕与动态避障。
\end{itemize}

\section{金字塔优先经验回放机制}
为解决长周期任务中的样本筛选问题，设计三级存储结构（PHM）：
L0（原始池）、L1（高 TD-Error 惊喜池）与 L2（成功轨迹池）。在稀疏奖励下，L2 样本是极高价值的“灯塔”。

\section{实验与结果分析}
\subsection{综合性能评估}
在 5 机围捕任务中，HGC-RL 的围捕成功率达到 92\%，平均完成时间比 MADDPG 减少 28\%，且碰撞率降低 60\%。
\subsection{消融实验与案例分析}
去掉金字塔经验池后，在长达 5000 步的长时序任务中，仅使用标准 Replay Buffer 的 Agent 在后期性能下降了 20\%（发生灾难性遗忘），而 PHM 保持了性能的单调上升。

\section{本章小结}
本章通过课程学习和金字塔经验池，实现了从单机技能到多机协同的高效策略迁移。
